\documentclass[a4paper,10pt]{article}

\usepackage{latexsym}
\usepackage[T1]{fontenc}
\usepackage[a4paper, top=0.3in, bottom=0.2in, left=0.45in, right=0.45in, headheight=0pt, headsep=2pt, footskip=4pt, includeheadfoot]{geometry}
\usepackage{titlesec}
\usepackage{marvosym}
\usepackage[usenames,dvipsnames]{color}
\usepackage{enumitem}
\usepackage{fancyhdr}
\usepackage[pdftex, colorlinks=true, urlcolor=blue, linkcolor=blue]{hyperref}

\pagestyle{fancy}
\fancyhf{}
\fancyfoot{}
\renewcommand{\headrulewidth}{0pt}
\renewcommand{\footrulewidth}{0pt}


\urlstyle{rm}

% pdfstartlink-based \href for guaranteed clickable links
\makeatletter
\newcommand{\hreflink}[2]{%
  \leavevmode
  \pdfstartlink attr{/Border [0 0 0]}%
    user{/Subtype /Link /A << /S /URI /URI (#1) >>}%
  {\color{blue}#2}%
  \pdfendlink
}
\let\oldhref\href
\renewcommand{\href}[2]{\hreflink{#1}{#2}}
\makeatother

\raggedbottom
\raggedright
\setlength{\tabcolsep}{0in}

\titleformat{\section}{
  \vspace{-11pt}\scshape\raggedright\large
}{}{0em}{}[\color{black}\titlerule \vspace{-7pt}]

\newcommand{\resumeItem}[2]{
  \item\small{\textbf{#1}{: #2 \vspace{-1pt}}}
}
% ATS-safe: \hfill keeps a single linear text flow (no tabular columns to jumble)
\newcommand{\resumeSubheading}[4]{
  \vspace{1pt}\item
    \textbf{#1}\hfill{#2}\\
    \textit{\small #3}\hfill\textit{\small #4}\vspace{-4pt}
}
\newcommand{\resumeSubItem}[2]{\resumeItem{#1}{#2}\vspace{-4pt}}
\newcommand{\resumeDesc}[1]{\item[]\small{#1}\vspace{-3pt}}
\renewcommand{\labelitemii}{-}
\newcommand{\resumeSubHeadingListStart}{\begin{itemize}[leftmargin=*,topsep=2pt,itemsep=0pt]}
\newcommand{\resumeSubHeadingListEnd}{\end{itemize}}
\newcommand{\resumeItemListStart}{\begin{itemize}[itemsep=1pt,topsep=1pt]}
\newcommand{\resumeItemListEnd}{\end{itemize}\vspace{-4pt}}

%----------AUTO-GENERATED PR COUNTS-----------------
% Everything between the PRSTATS markers is rewritten daily by
% scripts/update_contributions.py from the live GitHub API. Do not hand-edit:
% your changes will be overwritten. To use a count in the body, reference the
% macro (e.g. \prMerged{}) rather than typing a literal number.
% <!-- PRSTATS_START -->
\newcommand{\prMerged}{19}
\newcommand{\prOpen}{9}
\newcommand{\prRepos}{8}
\newcommand{\prMergedProjects}{PyTorch, vLLM, Meta Velox, Apache Spark, and Apache Gluten}
\newcommand{\prMergedPyTorch}{3}
\newcommand{\prMergedvLLM}{1}
\newcommand{\prMergedVelox}{5}
\newcommand{\prMergedSpark}{4}
\newcommand{\prMergedGluten}{6}
\newcommand{\prMergedDuckDB}{0}
% <!-- PRSTATS_END -->

\begin{document}

%----------HEADING-----------------
% ATS-safe: centered single-column contact block (linear text extraction)
\begin{center}
  {\LARGE \textbf{Brij Raj Kishore}}\\ \vspace{2pt}
  {\small Staff Software Engineer \;|\; \textbf{C++} \;|\; \textbf{Java} \;|\; \textbf{Python}}\\ \vspace{3pt}
  \small
  \href{mailto:brijraj31@gmail.com}{brijraj31@gmail.com} $|$ +91-7906555465 $|$ Noida, Delhi/NCR, India\\
  \href{https://www.linkedin.com/in/brijrajkishore/}{linkedin.com/in/brijrajkishore} $|$ \href{https://github.com/brijrajk}{github.com/brijrajk}
\end{center}

%-----------SUMMARY-----------------
\section{Professional Summary}
Systems software engineer (\textasciitilde{}6 yrs) in \textbf{high-performance C++, Java, and Python} for ML and data systems. Building a \textbf{PyTorch out-of-tree device backend for low-latency LLM inference}, on a foundation of query-engine internals (Velox, Gluten, Spark), \textbf{JVM} runtimes, and heterogeneous compute (FPGA, CUDA). \textbf{\prMerged{} merged PRs} across \prMergedProjects{}; \prOpen{} open across \prRepos{} repos.

%-----------WORK EXPERIENCE-----------------
\section{Work Experience}
\resumeSubHeadingListStart

  \resumeSubheading
    {ZettaBolt Technologies Pvt. Ltd}{Noida, India}
    {Senior Member of Technical Staff (Staff Software Engineer)}{December 2021 - Present}
    \resumeItemListStart
      \resumeItem{PyTorch OOT Device Backend for LLM Inference}
        {Building an out-of-tree \textbf{C++ device backend} in PyTorch (PrivateUse1) - operator dispatch, kernel registration, device memory, streams - against the \textbf{OpenReg} reference device ahead of custom silicon; validated via \textbf{OpInfo} and ran \textbf{Llama and BERT inference end-to-end}. Built a \textbf{CUDA-parity testing framework} + CI. Currently building the \textbf{Inductor/Triton codegen path} bringing \textbf{torch.compile} to the device, standalone and through \textbf{vLLM} serving.}
      \resumeItem{Agent Harness for Accelerator-Backend Development}
        {Built a domain-specific \textbf{agent harness} (Claude Code) for PyTorch backend work - \textbf{79 files / \textasciitilde{}7K lines}: 27 skill areas, \textbf{22 subagents}, and \textbf{11 hooks} enforcing citation-verified answers, generated-file protection, and blast-radius/test-impact analysis per edit. The \textbf{CUDA-parity} and \textbf{OpInfo} tooling above run on it.}
      \resumeItem{Spark Query Engine Acceleration via FPGA Offloading}
        {Achieved \textbf{5x speedup over vanilla Spark} and \textbf{2x over Apache Gluten} by integrating \textbf{Velox (C++)} with an FPGA SDK through a \textbf{Scala} bridge; designed the operator dispatch architecture for Scan, Filter, and Aggregation.}
      \resumeItem{AMD AOCL Integration on Meta Velox}
        {Delivered \textbf{5x runtime improvement} on AMD EPYC by integrating AMD Compiler and AOCL into Meta\textquotesingle{}s Velox/Gluten pipeline (\textbf{C++/Scala}); guided by \textbf{profiling} with uProf and Flamegraph.}
      \resumeItem{Graph Analytics: NVIDIA cuGraph \& TigerGraph}
        {Achieved \textbf{100x speedup} on PageRank using TigerGraph GSQL with \textbf{NVIDIA CUDA (C++)} and Python; cut pipeline runtime \textbf{40\%} and infra cost \textbf{24\%}.}
      \resumeItem{ZettaProf - Automated Spark Profiling \& Benchmarking Tool}
        {\textbf{Reduced Spark debugging time 60\%} by designing and shipping \textbf{ZettaProf} (Scala/Java), an automated profiler adopted \textbf{company-wide}; mentored engineers on its use.}
      \resumeItem{JVM-to-C++ Hive Offload \& Healthcare ETL on Microsoft Fabric}
        {Improved Hive query runtime \textbf{20\%} migrating execution kernels from the \textbf{JVM to a C++ backend} for heterogeneous CPU/FPGA dispatch; separately \textbf{cut debugging time \textasciitilde{}70\%} optimizing healthcare pipelines on \textbf{Spark (PySpark/Scala)}, \textbf{Microsoft Fabric}, and \textbf{Azure Data Factory} (join skew, AQE fallback).}
    \resumeItemListEnd

  \resumeSubheading
    {NXP Semiconductors}{Pune, India}
    {Software Engineer}{July 2021 - December 2021}
    \resumeDesc{Ported BT firmware to dual-antenna platform in \textbf{C}; debugged WiFi/BLE driver bugs via GDB; validated via \textbf{unit tests}.}

  \resumeSubheading
    {Tata Consultancy Services}{Bengaluru, India}
    {Assistant System Engineer}{Aug 2017 - Jan 2018}
    \resumeDesc{Developed \textbf{REST API} endpoints (TCS OmniStore POS); raised unit test coverage \textbf{62\% to 70\%} via JUnit and \textbf{CI/CD}.}

\resumeSubHeadingListEnd

%-----------OPEN SOURCE-----------------
\section{Open Source Contributions \textnormal{\small(\prMerged{} merged, \prOpen{} open across \prRepos{} upstream repos)}}
\resumeSubHeadingListStart
  \resumeSubItem{Merged}
    {\href{https://github.com/pytorch/pytorch}{\textbf{PyTorch}} (\prMergedPyTorch{}) - CUDA device-guard RAII fix on the throwing path, \texttt{LPPoolNd} validation, schema-inference errors; \href{https://github.com/facebookincubator/velox}{\textbf{Velox}} (\prMergedVelox{}) - TPC-DS map-allocation perf, Spark \texttt{transform\_values} registration, Parquet encoding; \href{https://github.com/apache/gluten}{\textbf{Gluten}} (\prMergedGluten{}) - \texttt{KeyGroupedPartitioning} \texttt{MatchError} fix, codec error messaging, test-suite fixes; \href{https://github.com/apache/spark}{\textbf{Spark}} (\prMergedSpark{}) - config/docs fixes; \href{https://github.com/vllm-project/vllm}{\textbf{vLLM}} (\prMergedvLLM{}) - Transformers v5 test gating.}
  \resumeSubItem{Under review}
    {PyTorch \texttt{fp16} CPU/GPU parity (\texttt{clamp}, \texttt{torch.where}); Gluten bloom-filter AQE-fallback corruption; Spark structured streaming (\texttt{GroupState}, Protobuf); DuckDB \texttt{*COLUMNS()} resolution.}
\resumeSubHeadingListEnd

%-----------TECHNICAL SKILLS-----------------
\section{Technical Skills}
\resumeSubHeadingListStart
  \resumeSubItem{ML / AI}{\hspace{1.2em}PyTorch (OOT device backend, PrivateUse1, OpenReg, dispatch, kernel registration, OpInfo, streams), \textbf{torch.compile} (Inductor, Triton), CUDA-parity testing, NVIDIA CUDA, cuGraph, vLLM}
  \resumeSubItem{Languages}{\hspace{0.5em}C++17, C, \textbf{Java}, Scala, Python, Bash}
  \resumeSubItem{Query \& Distributed}{\hspace{0.1em}Meta Velox, Apache Gluten, Spark SQL, Hive, DuckDB, Parquet/ORC, Hadoop, Kafka, JVM internals, Concurrency}
  \resumeSubItem{Hardware \& Tools}{FPGA (Xilinx/AMD), AMD EPYC / AOCL, GDB, Valgrind, uProf, Flamegraph, Docker, Linux, Git, CI/CD}
  \resumeSubItem{Cloud}{\hspace{2.4em}AWS, Microsoft Fabric, Azure Data Factory}
\resumeSubHeadingListEnd

%-----------EDUCATION-----------------
\section{Education}
\resumeSubHeadingListStart
  \resumeSubheading
    {National Institute of Technology (NIT Nagpur)}{Nagpur, India}
    {M.Tech - Computer Science \& Engineering;\quad CGPA: 8.04\,/\,10}{July 2019 - June 2021}
  \resumeSubheading
    {GL Bajaj Institute of Technology \& Management}{Greater Noida, India}
    {B.Tech - Computer Science \& Engineering}{July 2013 - June 2017}
\resumeSubHeadingListEnd

%-----------ACHIEVEMENTS-----------------
\section{Achievements}
\vspace{2pt}
\small{Author, \href{https://github.com/brijrajk/pytorch-internals-explained}{\textit{pytorch-internals-explained}} - public series on PyTorch internals (dispatch, autograd, strides/views, CUDA allocator) \quad|\quad GATE CSE 2019 - \textbf{97\textsuperscript{th} Percentile} \quad|\quad Google Code Jam 2020 - \textbf{6,448\,/\,44,434} \quad|\quad TCS CodeVita 2016 - \textbf{414 AIR}}

\end{document}
